%% ma: L12-B08-0045
%% lop: 12
%% chuong: 2
%% bai: 8
%% dang: 12.8.32
%% loai: TN4
%% muc: TH
%% dap_an: "A"
%% nguon: "LaTeX giáo viên gửi 21/10/2026 (M05, Vectơ và hệ trục tọa độ trong không gian, 107 câu), Phần I Câu 49"
%% kien_thuc: "Bài 8 (biểu thức tọa độ của tích vô hướng), ứng dụng tính công"
%% trang_thai: cho-duyet
%% ly_do: "Dạng mới đề xuất 12.8.32"
\begin{ex}
Một thiết bị kéo hàng tác dụng một lực $\overrightarrow{F}=(100;200;50)$ (đơn vị: N) làm vật dịch chuyển theo vectơ độ dời $\vec{d}=(4;3;-2)$ (đơn vị: m). Công sinh ra bởi lực $\overrightarrow{F}$ là
\choice{\True $900$ J}{$1100$ J}{$1000$ J}{$800$ J}
\loigiai{$A=\overrightarrow{F}\cdot\vec{d}=400+600-100=900$ (J).}
\end{ex}
